\chapter{Admissibility as a Logic}
\label{ch:admissibility}

The preceding two chapters relaxed classical logic in two directions
independently: extra truth values, and truth relativized to a state.
This chapter's claim is that a third, related idea — admissibility —
is not a metaphor borrowed from those two but a logic in its own
right, with its own notion of what follows from what. A history of
witnessed facts only grows; a constraint can narrow which future
continuations are admissible without rewriting anything already
witnessed. Both are ordinary logical claims once stated precisely,
and both are provable rather than merely evocative.

\section{Histories and Extension}

\begin{experiment}{Extension is a preorder}
\begin{lean}
structure History (alpha : Type) where
  witnessed : Set alpha

def Extends {alpha : Type} (earlier later : History alpha) : Prop :=
  earlier.witnessed ⊆ later.witnessed

theorem extends_refl {alpha : Type} (h : History alpha) :
    Extends h h := fun x hx => hx

theorem extends_trans {alpha : Type} {h1 h2 h3 : History alpha}
    (h12 : Extends h1 h2) (h23 : Extends h2 h3) :
    Extends h1 h3 := fun x hx => h23 (h12 hx)

theorem witnessed_facts_persist {alpha : Type}
    {earlier later : History alpha}
    (h : Extends earlier later) {fact : alpha}
    (w : fact ∈ earlier.witnessed) : fact ∈ later.witnessed :=
  h w
\end{lean}
Reflexivity and transitivity make \texttt{Extends} a preorder on
histories — the same structural fact that, read as an accessibility
relation, would license the transitive modal reasoning flagged at
the end of Chapter~\ref{ch:modal}.
\end{experiment}

\section{Constraint as Monotone Narrowing}

\begin{experiment}{Constraining never enlarges what is admissible}
\begin{lean}
def Constrain {alpha : Type} (admissible constraint : Set alpha) : Set alpha :=
  admissible ∩ constraint

theorem constrain_subset {alpha : Type}
    (admissible constraint : Set alpha) :
    Constrain admissible constraint ⊆ admissible := fun x hx => hx.1
\end{lean}
This is deliberately almost too simple to call a theorem — the
point is that it \emph{is} one: ``constraint can only narrow, never
expand, what a process may still do'' is not asserted rhetorically
anywhere in this framework, it is this one-line fact, checked the
same way every other claim in this book has been checked.
\end{experiment}

\section{A Consequence Relation for Admissibility}

A logic needs more than isolated facts about \texttt{Extends} and
\texttt{Constrain}; it needs a notion of what follows from what. Say
a fact $f$ is \emph{admissible} relative to a history $h$ and a
constraint $C$ when $f$ is both witnessed and permitted — exactly
the \texttt{Constrain} of Section~2, given a name.

\begin{experiment}{Two structural facts, honestly named}
\begin{lean}
def Admissible {alpha : Type} (h : History alpha) (C : Set alpha)
    (f : alpha) : Prop := f ∈ Constrain h.witnessed C

theorem admissible_of_witnessed_and_permitted {alpha : Type}
    (h : History alpha) (C : Set alpha) (f : alpha)
    (hw : f ∈ h.witnessed) (hc : f ∈ C) : Admissible h C f :=
  And.intro hw hc

theorem admissible_persists_under_extension {alpha : Type}
    {h h' : History alpha} (hext : Extends h h')
    {C : Set alpha} {f : alpha} (ha : Admissible h C f) :
    Admissible h' C f :=
  And.intro (hext ha.1) ha.2
\end{lean}
The first is admissibility's version of a reflexivity/assumption
rule: a fact that is directly witnessed and directly permitted is
admissible with no further argument needed. The second is a
persistence fact under history extension, built from exactly the
same accessibility-composition idea as
Chapter~\ref{ch:modal}'s \texttt{committed\_reaches\_recorded}: an
admissible fact stays admissible as the history grows, because
\texttt{hext} carries the witnessed half forward and the constraint
half never changed.
\end{experiment}

Both facts hold with the constraint $C$ held fixed. What makes
admissibility a genuinely different logic — not merely a restatement
of the classical or modal machinery already built — is what happens
when $C$ itself is allowed to change.

\begin{countermodel}{Admissibility is not monotone in the constraint}
\begin{lean}
def sampleHistory : History Nat where
  witnessed := {n | n = 7}

def openConstraint : Set Nat := {n | True}
def stricterConstraint : Set Nat := {n | n ≠ 7}

example : Admissible sampleHistory openConstraint 7 :=
  And.intro rfl trivial

example :
    ¬ Admissible sampleHistory (openConstraint ∩ stricterConstraint) 7 := by
  intro h
  exact h.2.2 rfl
\end{lean}
$7$ is admissible under the open constraint — it is witnessed, and
the open constraint permits everything. Narrowing the constraint by
intersecting in \texttt{stricterConstraint} does not add any new
witnessed facts and does not touch \texttt{Extends} at all, yet it
removes $7$'s admissibility outright. In classical logic, adding
premises can only add consequences, never remove them — that is
what monotonicity means. Here, adding a constraint can \emph{retract}
a previously admissible conclusion. This is not a defect to be
patched; it is the entire content of what ``constraint'' is supposed
to mean, and it is the specific respect in which admissibility is a
logic with its own structural profile rather than classical or
modal logic wearing new vocabulary.
\end{countermodel}

\section{Exercises}

\begin{exercise}
Prove that \texttt{Extends} \emph{is} antisymmetric for
\texttt{History} exactly as defined here: if two histories extend
each other in both directions, their witnessed sets are equal by
set extensionality, and since \texttt{witnessed} is the structure's
only field, equal witnessed sets make the histories themselves
equal.
\end{exercise}

\begin{exercise}
Antisymmetry above relied on \texttt{witnessed} being
\texttt{History}'s only field. Define an enriched structure,
\texttt{RichHistory}, adding a second field (a provenance tag, say,
of type \texttt{String}), with \texttt{Extends} defined to compare
only the \texttt{witnessed} components as before. Exhibit two
distinct \texttt{RichHistory} values — same witnessed set, different
provenance tags — that extend each other in both directions without
being equal, confirming that antisymmetry does not survive this
natural enrichment.
\end{exercise}
